\def\OutcomeAuditIntegrated{1}
\section{Buying the information that missing customers take away}
\label{closure30:outcome-audit}

A bank cannot recover the experience of customers who disappear merely by
fitting a more elaborate model to those who remain. It can sometimes change
the evidence. A sample of incomplete customer histories may be recovered from
archived servicing records, linked payments or carefully conducted follow-up.
The economic question is how much reliable information that extra work buys.

\ifdefined\OutcomeAuditIntegrated
The preceding analysis shows why resampling a fitted model does not correct
outcome-dependent dropout. This section changes the observation design while
retaining its three-decision servicing example.
\else
This example continues a generated servicing comparison in which dropout
depends on an unseen future shock. Resampling a fitted transition model can
reproduce that selection error and produce misleadingly narrow intervals.
\fi
Random recovery repairs the missing-information problem under a clear
condition: the histories actually recovered must form the specified random
sample. Random invitations are insufficient if successful recovery remains
selective. Even a correct estimate can be too imprecise to support a decision.
Here a separately fitted outcome model reduces interval width substantially,
but the additional pilot information is costly and a positive policy
improvement often remains statistically unresolved.

\subsection{A known design, applied to a sequential decision}
Following up a subsample with missing outcomes is usually called double
sampling. Levis and colleagues show how noninformative successful second-stage
recovery identifies the full-data distribution despite informative original
missingness. Their Proposition~1 treats a missing outcome; Appendix~A extends
the argument to general coarsening of a record. Their mean-effect estimators
also use regression to improve efficiency under explicit conditions
\citep{levis2023double}.

This changes the information without exempting the analysis from selection
assumptions. Levis and colleagues distinguish invitations from obtaining the
missing information, and discuss nonresponse and recall error. Sun and
colleagues develop a related construction for weighted quantile treatment
effects, with conditions appropriate to quantiles \citep{sun2025double}.
A mean-reward comparison does not thereby qualify a tail-loss estimate.
Averaging conditional quantile effects does not in general recover a
population quantile effect.

The present contribution is a worked composition: recover missing pieces of
a sequential policy-value score, derive the precision consequence, and test
it in the servicing law. The two-phase idea and sequential regression
assistance are inherited methods. This does not establish an efficient
estimator for every longitudinal missing-data model.

\subsection{What must be recovered}
There are three decisions, $t=0,1,2$. Each generated customer has a fixed
characteristic $X\sim N(0,1)$ and initial state
$L_0=0.4X+\varepsilon_0$. With independent standard-normal innovations,
\begin{align*}
L_{t+1}&=0.6L_t+0.25X-0.5A_t+\varepsilon_{t+1},\\
R_t&=10-2L_{t+1}^{2}-A_t+0.5A_tX.
\end{align*}
The historical probability of action is
\[
p_t=\operatorname{clip}\{\operatorname{expit}(0.3X-0.5L_t+0.2A_{t-1}),
                       0.2,0.8\},\qquad A_{-1}=0.
\]
These are known probabilities in a generated example. Clipping belongs to
the assignment law, rather than modifying an estimated bank propensity
score. A fixed policy acting when $L_t>0.5$ is compared with never acting.
Neither is selected using evaluation outcomes.

For history $H_t$ available before action $t$, define
\[
W_t(\pi)=\prod_{s=0}^{t}
 \frac{\mathbf1\{A_s=\pi(H_s)\}}{p_s(A_s\mid H_s)},
\qquad F(\pi)=\sum_{t=0}^{2}W_t(\pi)R_t.
\]
Sequential exchangeability, consistency and positive action probabilities
give $E\{F(\pi)\}=V(\pi)$, where
$V(\pi)=E_\pi(\sum_t R_t)$. The simulation satisfies these conditions by
construction. Recovering missing records cannot make them true in a bank.

Let $C_t$ indicate that reward $t$ and its required history were routinely
observed. After dropout, later $C_t$ are zero. Split the full score:
\[
K(\pi)=\sum_t C_tW_t(\pi)R_t,\qquad
M(\pi)=\sum_t(1-C_t)W_t(\pi)R_t,\qquad F=K+M.
\]
The missing contribution may depend on unobserved outcomes. Among incomplete
records, let $B$ indicate completed recovery of the missing history, with
known probability $q>0$, independently of that history and its original
observation pattern. Complete records need no recovery and have $M=0$.
Average the following score over all initially eligible customers:
\begin{equation}
Z(\pi)=K(\pi)+\frac{B}{q}M(\pi).
\label{closure30:audit-score}
\end{equation}
Each routinely observed contribution has weight one; each recovered missing
contribution represents $1/q$ similar incomplete histories. There is no
division by the questionable routine observation probability.

Condition on the full history and routine observation pattern.
$E(B\mid\text{history},C)=q$ for incomplete records, so
$E(Z\mid\text{history},C)=F$ and $E(Z)=V(\pi)$. This requires recovery of
the quantities entering $M$: actions, states, rewards and the information
needed for assignment probabilities. A terminal balance alone may not
supply them. Recovery must measure historical outcomes faithfully and must
not change them. A follow-up intervention that changes repayment would
answer a different question.

\subsection{Identification does not imply useful precision}
Under constant independent Bernoulli recovery,
\begin{equation}
\operatorname{Var}(Z)=\operatorname{Var}(F)
                   +(q^{-1}-1)E(M^2).
\label{closure30:audit-variance}
\end{equation}
Indeed $Z-F=(B/q-1)M$ has conditional mean zero and conditional variance
$(q^{-1}-1)M^2$, so its covariance with $F$ is zero. For policy improvement,
use paired differences $F_\Delta=F(\pi)-F(0)$ and
$M_\Delta=M(\pi)-M(0)$. Evaluating both policies on the same customers
retains this pairing.

Low recovery probabilities magnify rare contributions. Incomplete customers
who match several consecutive policy actions can receive large weights.
More recovery reduces the extra variance; recovering every history still
leaves the variance of the full-history policy score.

The first routine observation mechanism continues with probability
$g_t=0.65+0.30\operatorname{expit}(0.3L_t-0.2A_t)$, independently of the
next innovation conditional on history and action. The adverse mechanism
uses $\operatorname{expit}\{\operatorname{logit}(g_t)+1.5\varepsilon_{t+1}\}$.
Weighting by the original $g_t$ in that case assumes away selection by the
unseen shock.

Each mechanism has 500 independent replications with 2,000 customers.
Recovery rates are 5\%, 10\%, 25\%, 50\% and 100\% of incomplete records,
with shared recovery draws across rates within a replication. Estimators
receive full histories only where routine observation or designated recovery
permits. The reference policy improvement is 0.75258 reward units, with
numerical Monte Carlo standard error 0.00432, from the retained
million-trajectory intervention calculation. Reward units are illustrative,
not dollars.

\begin{table}[htbp]
\centering\small
\begin{tabular}{lrrrr}
\toprule
Adverse dropout & Bias & Coverage & Width & Follow-ups\\
\midrule
Assumed routine $g_t$ & 0.755 & 0.908 & 4.506 & 0\\
Random recovery, 5\% & -0.129 & 0.944 & 11.809 & 62\\
Random recovery, 25\% & -0.051 & 0.948 & 6.047 & 308\\
Random recovery, 50\% & 0.016 & 0.930 & 4.772 & 618\\
Random recovery, 100\% & 0.029 & 0.942 & 3.933 & 1,238\\
\bottomrule
\end{tabular}
\caption{Fixed-policy improvement under outcome-dependent dropout;
500 replications, 2,000 evaluation customers. Width is the mean width of a
normal 95\% interval. Follow-ups are completed missing-history recoveries.
No pilot is used. Coverage is against the numerical reference, not a
finite-sample guarantee.}
\end{table}

The estimated bias at 25\% recovery is small relative to its Monte Carlo
standard error, 0.070. Yet mean interval width is about eight times the
improvement. Even full recovery yields a positive lower bound in only
11.6\% of replications. Under conditionally independent routine dropout,
25\% recovery gives width 5.679 and coverage 93.0\%. These results do not
establish uniform attainment of nominal coverage.

A third study makes actual recovery depend on the unseen first-missed
innovation, using
$\operatorname{expit}\{\operatorname{logit}(0.25)+0.8\varepsilon_{\rm miss}\}$,
but analyses it as constant 25\% recovery. Across 500 replications, upward
bias is 0.278 with Monte Carlo standard error 0.063. Coverage nevertheless
appears to be 95.2\%, because mean width is 5.219. A nominal-looking
coverage fraction alone does not validate the design or exclude bias.

\subsection{Use prediction to economise on information}
Repeated weighting expends information on action randomness as well as
reward uncertainty. A model of future reward can remove part of this
variation. The sequential doubly robust construction of
\citet{jiang2016}, equation~(10) and Theorem~1, makes this distinction
precise. Here the model is fitted on a separate pilot and action
probabilities remain known. Unbiasedness does not require the fitted
outcome model to be correct.

For fixed pilot regressions $\widehat Q_t(H_t,A_t)$, put
$\widehat V_t(H_t)=\widehat Q_t\{H_t,\pi(H_t)\}$ and
$\widehat V_3=0$. The full-history score is
\begin{equation}
D(\pi)=\widehat V_0(H_0)+
 \sum_{t=0}^{2}W_t(\pi)
 \{R_t+\widehat V_{t+1}(H_{t+1})-\widehat Q_t(H_t,A_t)\}.
\label{closure30:sequential-score}
\end{equation}
Condition on the pilot. With $W_{-1}=1$, random assignment gives
\[
E\{W_t\widehat Q_t(H_t,A_t)\}
  =E\{W_{t-1}\widehat V_t(H_t)\}.
\]
Model terms cancel in expectation, leaving
$E(D\mid\text{pilot})=V(\pi)$. This telescoping argument covers the present
reward and transition dependence without importing a variance formula under
unstated assumptions. Setting every $\widehat Q_t$ to zero returns the
reward-weighted score.

Apply recovery to the residual contributions. Initial $\widehat V_0$ is
observed for everyone and belongs in $K_D$; split the residual sum into
routinely observed $K_D$ and missing $M_D$. Then $Z_D=K_D+B M_D/q$ has
conditional mean $D$ and
\[
\operatorname{Var}(Z_D\mid\text{pilot})
=\operatorname{Var}(D\mid\text{pilot})
 +(q^{-1}-1)E(M_D^2\mid\text{pilot}).
\]
Good predictions can reduce both terms. A wrong but fixed square-integrable
model does not invalidate the mean identity when the probabilities are
correct. Reusing evaluation outcomes for fitting would require a different
argument, such as a justified sample-splitting procedure.

\subsection{The precision improvement and its price}
A second experiment uses 5,000 independent pilot customers per
replication and random 25\% recovery of incomplete pilot histories.
Working backward, separate regressions for each policy use the fixed basis
$(1,L,X,A,L^2,X^2,LX,AL,AX)$. Routine transitions receive weight one and
recovered transitions weight four. Another 2,000 customers evaluate the
scores. There are 500 replications per observation mechanism, with fresh
seeds. The model basis is fixed before final evaluation.

\begin{table}[htbp]
\centering\small
\begin{tabular}{lrrrr}
\toprule
Adverse dropout & Bias & Coverage & Width & Positive lower\\
\midrule
Assumed $g_t$, assisted & 0.926 & 0.590 & 2.195 & 84.4\%\\
5\% recovery, assisted & 0.082 & 0.958 & 4.770 & 15.8\%\\
25\%, reward weighting & -0.079 & 0.962 & 6.094 & 4.4\%\\
25\%, constant model & -0.025 & 0.968 & 3.229 & 14.4\%\\
25\%, quadratic model & -0.003 & 0.964 & 2.773 & 24.0\%\\
50\%, quadratic model & 0.016 & 0.950 & 2.181 & 33.2\%\\
Full, quadratic model & 0.001 & 0.952 & 1.827 & 41.4\%\\
\bottomrule
\end{tabular}
\caption{Fixed-policy improvement with an independent pilot. Rows share
evaluation samples and recovery draws. Assisted and constant-model rows
also use the 5,000-customer pilot. Positive lower denotes a positive
interval endpoint, not adoption. Reward weighting needs no pilot.}
\end{table}

At 25\% evaluation recovery, the quadratic model reduces width from 6.094
to 2.773, or 54.5\%. The paired difference is $-3.320$, with Monte Carlo
standard error 0.038. Bias is $-0.0026$ with Monte Carlo standard error
0.0323. Coverage of 96.4\% has Monte Carlo standard error 0.83 percentage
points. Under conditionally independent routine dropout, coverage is 95.2\%
and width 2.607. The crude constant model also remains centred within
simulation error, while its wider intervals show the cost of weaker
prediction.

The pilot is additional information. Under adverse dropout it needs about
775 completed recoveries on average, besides roughly 309 evaluation
recoveries at 25\%. The unassisted comparator needs only evaluation
recoveries. This is not an equal-total-budget comparison. Reusing a pilot
across policies might spread its cost, but that benefit must be measured.

\begin{figure}[htbp]
\centering
\includegraphics[width=\textwidth]{twelve-hour-closure/2026-10-29/outcome-figures/precision.pdf}
\caption{More recovery narrows intervals, but positive evidence for a modest
gain remains scarce. Outcome-dependent dropout; 500 replications per study.
The assisted curve needs about 775 extra pilot recoveries, excluded from the
horizontal evaluation count. Curves use separate studies; the
paired 25\% comparison is in the text.}
\end{figure}

The false routine observation model with the assisted score produces a more
dangerous result: bias 0.926 and coverage 59.0\%, while 84.4\% of intervals
have a positive lower endpoint. Prediction sharpens a selected view of the
customers without correcting the observation mechanism.

\begin{figure}[htbp]
\centering
\includegraphics[width=\textwidth]{twelve-hour-closure/2026-10-29/outcome-figures/bias-coverage.pdf}
\caption{The observation design controls bias; regression controls part of
the variance. Bars are ordinary 95\% Monte Carlo intervals across 500
replications, not simultaneous qualification bounds. Reference uncertainty
is separate. Full history is a diagnostic benchmark, not bank data.}
\end{figure}

\subsection{A collection budget is a design choice}
Let $r$ be the incomplete-record probability, $c_0$ the cost of an initial
record and $c_A$ the additional cost of a recovered history. Expected
evaluation cost is $n(c_0+rqc_A)$, plus pilot cost. For a fixed score let
$V_F=\operatorname{Var}(F_\Delta)$ and $B_M=E(M_\Delta^2)$. Estimator
variance is $(V_F-B_M+B_M/q)/n$.

With fixed $n$, more recovery always reduces the missing-information term.
With a fixed monetary budget, it can mean fewer initial customers. If
$A=V_F-B_M>0$, differentiating
$(A+B_M/q)(c_0+rqc_A)$ gives the unconstrained optimum
\[
q^*=\sqrt{\frac{B_Mc_0}{A r c_A}},
\]
restricted to $[q_{\min},1]$. If $A\leq0$, that expression decreases with
$q$ and favours $q=1$ within this cost model. This assumes constant costs
and a fixed trained score, not an optimal joint pilot/evaluation design.
No bank prices are supplied, so no numerical recovery fraction is recommended.

Another route should be separated from the tested regression. For information
$O$ available before recovery and a separately trained predictor $m(O)$ of
the missing contribution,
\[
Z_m=K+m(O)+\frac{B}{q(O)}\{M-m(O)\}
\]
is unbiased under the known conditional recovery design even when $m$ is
wrong. Set $m=0$ for complete records, which need no recovery; the design
and allocation calculations concern incomplete records. Its extra variance is
$E[\{q(O)^{-1}-1\}\{M-m(O)\}^2]$.
The best squared-error predictor is $E(M\mid O)$. With known residual
second moment $a(O)=E[\{M-m(O)\}^2\mid O]$ and cost $c(O)>0$, minimising
this term at a fixed expected recovery budget yields
$q(O)\propto\sqrt{a(O)/c(O)}$, subject to positivity and the ceiling one.
This is a derivation, not a tested allocation algorithm. A separate pilot
must estimate the quantities; fresh data must test the design. Unseen
future outcomes cannot determine which customer is recovered.

\subsection{What this changes in the proposal}
A prototype can now be proposed around a specific design: accept declared
recovery probabilities, retain the initial customer denominator, estimate
paired policy improvements, and display the extra information burden and
remaining uncertainty. It should distinguish dropout weights from recovery
weights and show when an outcome model improves only precision. Where
recoverable histories do not exist, sensitivity or partial-identification
analysis is the appropriate next step.

These results concern independent generated customers, known action
probabilities and fixed policies. They do not settle calendar dependence,
unmeasured treatment confounding, erroneous or partially recovered histories,
estimated treatment probabilities, or uncertainty after policy learning.
Normal intervals are assessed rather than declared exact. The simulation
keeps every unaudited missing value unavailable to fitting and estimation.
The generated evidence does not establish that real histories can be
recovered accurately, or that the collection expense is worthwhile.

The remaining question is whether a feasible collection design can deliver
accurate histories in sufficient numbers to separate the economic
alternatives at an acceptable total cost.
